%======================================================================
\section{Tails and moments}
\label{sec:tails}

Extrapolation is where smile models die, so the tails deserve their own
section --- both for what the model guarantees (exponential tails,
Lee-consistent wings, in closed form) and for what it does not (the fitted
tail parameters are priors, weakly identified by any realistic strip; the
section ends on that).

\subsection{Endpoint speeds are tail scales}

\begin{proposition}[Exponential tails]
\label{prop:tails}
For every $\theta\in\Theta_N$ there are constants $K_\pm>0$ with
\begin{equation}
\Prob(X>x)\sim K_+\,e^{-x/\lambda_+}\ \ (x\to+\infty),
\qquad
\Prob(X<x)\sim K_-\,e^{\,x/\lambda_-}\ \ (x\to-\infty),
\label{eq:exptails}
\end{equation}
where $\sim$ denotes asymptotic equivalence (ratio $\to1$).  In return
terms: $\Prob(Y>y)\sim K_+y^{-1/\lambda_+}$ and
$\Prob(Y<y)\sim K_-y^{1/\lambda_-}$ --- exact power tails, hence regularly
varying with trivial slowly varying part.
\end{proposition}

\begin{proof}
Take the right tail; the left is a mirror image.  By
\eqref{eq:asymptotes}, $x(z)=(m+b_{+})+\lambda_+z+\varepsilon(z)$ with
$\varepsilon(z)=O(e^{-z})$ --- note the \emph{shifted} constant
$m+b_{+}$, since $x=m+\bar{x}$.  Now $\Prob(X>x)=1-\Lambda(z(x))$ where
$z(x)$ is the inverse transport.  Invert the asymptote explicitly:
evaluating the transport identity at $z=z(x)$,
\[
x=(m+b_{+})+\lambda_+z(x)+\varepsilon\big(z(x)\big)
\qquad\Longrightarrow\qquad
z(x)=\frac{x-(m+b_{+})}{\lambda_+}-\frac{\varepsilon(z(x))}{\lambda_+}
=\frac{x-(m+b_{+})}{\lambda_+}+O\big(e^{-z(x)}\big),
\]
the error keeping its exponential form because it is the same
$\varepsilon$, merely moved to the other side of the equation and
divided by the constant $\lambda_+$.  With
$1-\Lambda(z)=e^{-z}/(1+e^{-z})\sim e^{-z}$,
\[
\Prob(X>x)\sim e^{-z(x)}
= e^{-(x-m-b_{+})/\lambda_+}\,e^{O(e^{-z(x)})}
\sim e^{(m+b_{+})/\lambda_+}\,e^{-x/\lambda_+},
\]
which is \eqref{eq:exptails} with $K_+=e^{(m+b_{+})/\lambda_+}$.  The
return statements substitute $x=\log y$.
\end{proof}

The endpoint speeds are therefore genuine tail \emph{scales}: $\lambda_+$
is the number of log-return units the right tail needs to thin by a factor
$e$.  This is the sense in which the two leading coordinates of the model
are readable.  The tails are heavier than log-normal --- exponential in
$X$, power-law in $Y$ --- which is the empirically comfortable side to err
on for equity underliers, and the reason the smiles of
\cref{fig:transport} turn convexly upward in the wings.

\subsection{The moment strip}

Moments follow from one change of variables, done once, slowly, because the
entire admissibility story drops out of it.  Any expectation is a rank
integral, and in log-odds the rank measure carries the tent-shaped weight
$\rho(z)\asymp e^{-|z|}$:
\begin{equation}
\E\big[h(X)\big]
=\int_0^1 h\big(Q(u)\big)\,\dd u
=\int_{-\infty}^{\infty}h\big(x(z)\big)\,\rho(z)\,\dd z .
\label{eq:rankintegral}
\end{equation}

\begin{proposition}[Moment strip]
\label{prop:strip}
For every $\theta\in\Theta_N$,
\begin{equation}
\E\big[e^{rX}\big]<\infty
\qquad\Longleftrightarrow\qquad
-\frac{1}{\lambda_-}\;<\;r\;<\;\frac{1}{\lambda_+}\,.
\label{eq:strip}
\end{equation}
Equivalently, in return moments: the last finite positive moment order of
$Y$ beyond the first is $r_+^{*}=\sup\{r\ge0:\E[Y^{1+r}]<\infty\}
=1/\lambda_+-1$, and the last finite inverse moment order is
$r_-^{*}=\sup\{r\ge0:\E[Y^{-r}]<\infty\}=1/\lambda_-$.
\end{proposition}

\begin{proof}
Apply \eqref{eq:rankintegral} with $h(x)=e^{rx}$.  On the right,
$e^{rx(z)}$ grows like $e^{r\lambda_+z}$ (by the asymptotics in the proof
of \cref{prop:tails}) against the weight $\rho(z)\sim e^{-z}$: the
integrand behaves like $e^{(r\lambda_+-1)z}$, integrable at $+\infty$ iff
$r\lambda_+<1$; at equality the integrand tends to a positive constant
times $1$ and diverges.  On the left, $e^{rx(z)}\asymp e^{r\lambda_-z}$
against $\rho(z)\sim e^{z}$: integrable at $-\infty$ iff $r\lambda_->-1$.
Both endpoint failures are exact ties between the exponential growth of
the integrand and the exponential decay of the rank weight.  The return
statements are the substitutions $\E[Y^{1+r}]=\E[e^{(1+r)X}]$ (finite iff
$1+r<1/\lambda_+$) and $\E[Y^{-r}]=\E[e^{-rX}]$ (finite iff
$r<1/\lambda_-$).
\end{proof}

One point of the strip is load-bearing: the martingale integral is the
case $r=1$, and $1<1/\lambda_+$ is exactly the wall $\lambda_+<1$ of
\cref{prop:wall} --- the same fact seen for the third time, now as ``the
forward is the first moment, and the first moment must be finite''.  No
left-hand condition arises because $1>-1/\lambda_-$ always.

\subsection{Lee's formula, and the closed speed-to-slope maps}

Lee's moment formula \citep{Lee2004} is the model-free bridge from moments
to wing geometry.  Writing $\beta_R=\limsup_{k\to+\infty}w(k)/k$ and
$\beta_L=\limsup_{k\to-\infty}w(k)/|k|$ for the asymptotic slopes of the
total-variance smile, Lee proved
\begin{equation}
\beta_R=\Psi(r_+^{*}),
\qquad
\beta_L=\Psi(r_-^{*}),
\qquad
\Psi(r)=2-4\big(\sqrt{r^{2}+r}-r\big)\;\in\;[0,2],
\label{eq:lee}
\end{equation}
with $r_\pm^{*}$ the critical moment orders of \cref{prop:strip} (and
$\Psi(\infty)=0$, $\Psi(0)=2$: the fewer finite moments, the steeper the
admissible wing, capped at slope $2$).  For LQD the critical orders are
explicit, so the composition collapses to closed form.

\begin{proposition}[Speed-to-slope maps]
\label{prop:leeclosed}
For an LQD slice with endpoint speeds $\lambda_\pm$,
\begin{equation}
\beta_R=\Psi\Big(\frac{1}{\lambda_+}-1\Big)
=\frac{2\lambda_+}{\big(1+\sqrt{1-\lambda_+}\,\big)^{2}},
\qquad
\beta_L=\Psi\Big(\frac{1}{\lambda_-}\Big)
=\frac{2\lambda_-}{\big(1+\sqrt{1+\lambda_-}\,\big)^{2}} .
\label{eq:leeclosed}
\end{equation}
Both maps are increasing; $\beta_R\to2$ as $\lambda_+\to1$ (the Lee
ceiling is reached exactly at the wall) and $\beta_\pm\approx\lambda_\pm/2$
for small speeds.  Moreover, the $\limsup$s in \eqref{eq:lee} are genuine
limits for LQD slices.
\end{proposition}

\begin{proof}
Right wing: put $r=1/\lambda_+-1$, so $r^2+r=r(r+1)=(1-\lambda_+)/\lambda_+^2$
and $\sqrt{r^2+r}=\sqrt{1-\lambda_+}/\lambda_+$.  With
$s=\sqrt{1-\lambda_+}$,
\[
\Psi(r)=2-\frac{4}{\lambda_+}\big(s-s^{2}\big)
=2-\frac{4s(1-s)}{1-s^{2}}
=2-\frac{4s}{1+s}
=\frac{2(1-s)}{1+s}
=\frac{2(1-s^{2})}{(1+s)^{2}}
=\frac{2\lambda_+}{(1+s)^{2}} .
\]
Left wing: put $r=1/\lambda_-$, so $\sqrt{r^2+r}=\sqrt{1+\lambda_-}/\lambda_-$,
and with $t=\sqrt{1+\lambda_-}$,
\[
\Psi(r)=2-\frac{4}{\lambda_-}(t-1)
=2-\frac{4(t-1)}{t^{2}-1}
=2-\frac{4}{1+t}
=\frac{2(t-1)}{1+t}
=\frac{2(t^{2}-1)}{(1+t)^{2}}
=\frac{2\lambda_-}{(1+t)^{2}} .
\]
Monotonicity and the endpoint values are read off the final expressions.
The small-speed expansions follow from $s,t\to1$.  For the limit claim:
Lee's theorem gives only $\limsup$s in general, but Benaim and Friz
\citep{BenaimFriz2009} showed that when the return tail is regularly
varying --- as \cref{prop:tails} establishes here, with exact power tails
--- the wing slope converges as an ordinary limit.
\end{proof}

The endpoint coefficients therefore do double duty: they regularize the
extrapolation \emph{and} identify the entire moment budget and both
asymptotic wing slopes, in closed form, from two fitted numbers.
\Cref{fig:tails} draws the three-link chain --- endpoint speed to last
finite moment to Lee slope --- with the fitted values of two real
underliers marked on each link: for the SPY December 2026 slice of
\cref{sec:examples}, $\lambda_-=\MacSpyLamL$, $\lambda_+=\MacSpyLamR$,
$\beta_L=\MacSpyBetaL$, $\beta_R=\MacSpyBetaR$; for the long NVDA node
(December 2027, the $1.37$-year slice of \cref{fig:nvda_nodes}),
$\lambda_-=\MacNvdaLamL$, $\lambda_+=\MacNvdaLamR$,
$\beta_L=\MacNvdaBetaL$, $\beta_R=\MacNvdaBetaR$.  Single-name tails are
visibly heavier than index tails, as they should be --- and the NVDA
left scale is a live illustration of the wall's asymmetry:
$\lambda_->1$ is perfectly admissible (only the \emph{right} scale is
walled, \cref{prop:wall}); its price is a last finite inverse-moment
order $1/\lambda_-=\MacNvdaLongMomentLeft$ below one, i.e.\ even
$\E[1/Y]$ diverges, while the forward stays finite.

\begin{figure}[htbp]
\centering
\paperfig{\textwidth}{fig_tails.pdf}
\caption{The tail chain in three panels.  Panel A: endpoint speed
$\lambda$ to last finite moment order, $1/\lambda_-$ on the left tail and
$1/\lambda_+-1$ on the right \eqref{eq:strip}; the finite-forward wall
appears as the vertical asymptote of the right-tail map at
$\lambda_+=1$, where the right tail's moment budget is exhausted --- the
wall belongs to the right scale only, and the marked NVDA left scale
$\lambda_->1$ is legal.  Panel B: the Lee map $\Psi$ \eqref{eq:lee}.
Panel C: the composed speed-to-slope maps \eqref{eq:leeclosed}, with the
model-free ceiling $\beta=2$ dotted.  Markers on all three panels: the
fitted SPY December 2026 and NVDA December 2027 values from the frozen
snapshot --- equity fits use a small fraction of the admissible range,
sitting far below the ceiling.}
\label{fig:tails}
\end{figure}

\subsection{The honesty clause: the limit is not the wing}

Two cautions convert this section from advertisement to information.

First, the asymptotic slope is not a quote at any tradable strike ---
and at tradable strikes it can err in \emph{either direction}.  Near the
money the effective variance slope $w(k)/|k|$ is dominated by $w_0/|k|$
and diverges; far out it tends to $\beta$; but in between its path is
governed by the wing's intercept --- writing the wing as
$w(k)\approx\beta|k|+\text{const}$, the effective slope approaches the
limit from above when the constant is positive and from \emph{below}
when it is negative, and nothing pins the sign.  Both cases occur on one
real slice.  \Cref{fig:lee} plots the fitted SPY December node: on the
call wing the effective slope sits above its limit, ratio
\MacLeeRatioTenDeltaCall\ at the $10\Delta$ strike, still
\MacLeeRatioOneDeltaCall\ at $1\Delta$ --- an overshoot; on the put wing
it \emph{undershoots}, ratio \MacLeeRatioTenDeltaPut\ at $10\Delta$ and
\MacLeeRatioOneDeltaPut\ at $1\Delta$, both below one.  The curves in
the figure are truncated where the option price falls below $10^{-13}$
of forward --- beyond that, implied variance is Black-inversion noise
--- so the asymptote has not been approached at any strike a computer,
let alone a market, can quote.  The conclusion is therefore stronger
than a one-sided warning: a desk consuming $\beta_\pm$ should read them
as extrapolation boundary conditions and never price a wing off them
directly, because at ten-delta the limit overstates one wing of this
very slice and understates the other.

Second --- and this is the information clause of the paper's
validity/information/convention triad --- the fitted $\lambda_\pm$ are
\emph{priors dressed as parameters}.  No finite quoted strip determines an
asymptotic decay rate; what a fit reports as $\lambda_+$ is what the
quotes, the basis order, the ridge, the vega weighting, and the endpoint
coupling \emph{jointly} selected.  The model's contribution is to make
that selection explicit, low-dimensional, and overridable (a desk can pin
$\lambda_\pm$ through the endpoint chart of \cref{sec:computation}), not
to conjure tail information the market did not supply.
\Cref{sec:calibration} returns to identifiability with the fitting
machinery on the table.

\begin{figure}[htbp]
\centering
\paperfig{0.9\textwidth}{fig_lee.pdf}
\caption{The limit is not the wing --- in either direction.  Effective
total-variance slope $w(k)/|k|$ against $|k|$ (log scale) for the fitted
SPY December 2026 slice, both wings; the Lee limits $\beta_\pm$ of
\eqref{eq:leeclosed} dotted; diamonds mark the $10\Delta$ and $1\Delta$
strikes.  The call wing (right panel) sits \emph{above} its limit ---
ratios \MacLeeRatioTenDeltaCall\ at $10\Delta$, \MacLeeRatioOneDeltaCall\
at $1\Delta$; the put wing (left panel) dips \emph{below} its limit and
approaches it from underneath --- ratios \MacLeeRatioTenDeltaPut\ and
\MacLeeRatioOneDeltaPut.  Curves are truncated where the option price
drops below $10^{-13}$ of forward, beyond which implied variance is
inversion noise.  The asymptotic slope is a boundary condition on
extrapolation, not a quote at a tradable strike.}
\label{fig:lee}
\end{figure}
